\documentclass{article}

\usepackage[english]{babel}
\usepackage[letterpaper,top=2cm,bottom=2cm,left=3cm,right=3cm,marginparwidth=1.75cm]{geometry}
\usepackage{amsmath}
\usepackage{graphicx}
\usepackage[colorlinks=true, allcolors=blue]{hyperref}
% Scaffold mode: sections contain argument outlines as rendered bullet lists.
% Replace each itemize block with prose when ready to write.

\newcommand{\needpaper}{\textbf{[NEED PAPER]}}

\title{A Framework for Answer-Structure Prompting to Improve
Generalization and Accuracy of Few-Parameter LLMs}
\author{Eric Redman}

\begin{document}
\maketitle

% ============================================================
\begin{abstract}
\begin{itemize}
  \item Small LLMs are capable but underutilized due to unguided generation
  \item Existing prompting techniques address \textit{what} the model knows, not \textit{how} it structures its approach
  \item We propose a two-component pipeline: encoder classifier $\rightarrow$ few-shot scaffold $\rightarrow$ small LLM
  \item Central hypothesis: supplying a small LLM with an answer structure matched to the question improves its accuracy, narrowing the gap with larger models
  \item Enabling assumption: a finite taxonomy of answer structure types exists and is learnable by a classifier
  \item Experimental scope: multiple choice QA; evaluated across model sizes and configurations
  \item Result preview: \textbf{TBD after experiments}
\end{itemize}
\end{abstract}

% ============================================================
\section{Introduction}
% ============================================================

\subsection{Hook}

\begin{itemize}
  \item Large language models have transformed how we solve problems across domains, but their capabilities are bound to scale --- the strongest models are enormous, costly to run, and out of reach for most \cite{eloundou2023gpts, strubell2019energy}
  \item Small, open models are accessible to nearly anyone, yet they trail their larger counterparts on general tasks, leaving much of their promise unrealized \cite{smalllmsready2024, slm_survey2024}
  \item The usual explanation is scale --- with fewer parameters, they retain less of what they learn \cite{kaplan2020scaling, costofdownscaling2023}
  \item Yet their language ability is largely preserved \cite{linguistic_competence_scale2025} --- so is limited capacity really the whole story?
\end{itemize}

\subsection{The Problem Statement}

\begin{itemize}
  \item Large language models are a transformative technology --- reshaping how humans interact with information, create, and solve problems \cite{eloundou2023gpts, llm_critical_societal2024, llmoverview2023}
  \item But they are not open --- the most capable models are large, closed-source, and accessible only through proprietary APIs \cite{liesenfeld2024closed}
  \item Concentrated control limits who can contribute new ideas and perspectives to the field \cite{kapoor2024societal}
  \item It also determines what those models learn and how they act --- with no external input on training decisions \cite{bommasani2023fmti, liesenfeld2024closed}
  \item Closed-source models cannot be meaningfully audited, understood, or regulated by external parties \cite{bommasani2023fmti}
  \item And even then, many models marketed as open-source aren't truly open --- they provide weights for inference but not the data or infrastructure needed to train them \cite{liesenfeld2024openwashing, liesenfeld2024closed}
  \item Large models carry disproportionate environmental cost in training and inference \cite{strubell2019energy, patterson2021carbon}
  \item Only well-resourced organizations can build or deploy them at scale \cite{strubell2019energy, slm_survey2024}
  \item These constraints are inherent to large models and cannot be resolved through openness alone --- unfortunately, smaller models carry their own limitations
  \item Small open-source models exist but underperform due to limited parametric capacity \cite{modelsizematters2024, kaplan2020scaling}
  \item Small models' lower parameter counts impose a hard cap on retained knowledge and connections \cite{kaplan2020scaling}
  \item Small models cannot encode relational structure for every concept they were trained on --- they retain fragments but lack the capacity to organize all of them \cite{costofdownscaling2023}
  \item This limited capacity manifests as brittleness outside the training distribution and degraded performance on multi-step reasoning and complex real-world queries \cite{costofdownscaling2023}
  \item Despite these limitations, both large and small models share a common ability to understand language \cite{linguistic_competence_scale2025, emergent_reducedscale2024, costofdownscaling2023}
  % NOTE: linguistic_competence_scale2025 finds model size tied to memorization more than core language ability; emergent_reducedscale2024 shows small generative models retain emergent abilities; costofdownscaling2023 shows 60-70% reduction preserves in-context processing. Together these strongly support this claim for decoder LLMs specifically.
  \item The deficit in small models is not linguistic competence but associative depth --- they understand language but lack the structural scaffolding to organize what they do retain \cite{linguistic_competence_scale2025, costofdownscaling2023}
  \item And unlike their larger counterparts, small models are accessible enough for individual researchers to train and modify \cite{slm_survey2024}
  \item Small models can be trained effectively on specific domains, but this specialization comes at the cost of general capability --- a model trained for one domain cannot easily transfer to another \cite{domain_expert_smalllm2024, catastrophic_forgetting2023}
  \item Even so, small models still underperform on general tasks --- numerous approaches exist to address this, including prompt engineering, retrieval augmentation, and architectural specialization; we review them in detail in Section 2
  \item These approaches address knowledge, reasoning paths, or weight adaptation --- expanding what the model knows through retrieval, exploring multiple reasoning paths, or reshaping the weights through specific training --- but none provide a framework for how to structure the response
  \item We hypothesize that external answer structure may allow small models to approach the effectiveness of their larger counterparts
  \item \textbf{Research question:} can external answer structure compensate for limited capacity, making small open-source models effective enough to narrow the gap?
\end{itemize}

\subsection{The Purpose of the Paper}

\begin{itemize}
  \item This paper proposes and evaluates a framework that supplies small language models with external answer structure at inference time, testing whether it can narrow the gap identified in the problem statement
  \item \textbf{Central hypothesis:} conditioning a small LLM on an answer structure matched to the question improves its accuracy enough to narrow the gap with larger models. This is the claim the experiments are designed to test
  \item We propose a two-component inference-time pipeline requiring no weight updates to the generator:
    \begin{itemize}
      \item \textbf{Component 1:} an encoder-based classifier (BERT-style) maps a question $\rightarrow$ its answer structure type \cite{devlin2019bert}
      \item \textbf{Component 2:} any small LLM generates the answer conditioned on the retrieved few-shot scaffold corresponding to that structure type
    \end{itemize}
  \item The generator can be any small model that understands English \cite{smalllmsready2024, slm_survey2024}
  \item \textbf{Core mechanism:} the scaffold supplies the \textit{dimensions that matter} for a given question, not the answer itself --- the model fills them using its in-context learning ability, which withstands down-scaling even as parametric fact recall declines, drawing on the relational fragments it retains or retrieves \cite{costofdownscaling2023}
  \item The scaffold narrows the probability space \textit{before} generation begins, turning an unguided generation problem into a constrained one --- few in-context examples sharply reduce the model's uncertainty over the latent task \cite{xie2022icl}
  \item \textbf{Operating premises (assumed, not proven):}
    \begin{itemize}
      \item a finite set of answer structure types is sufficient to cover a given task --- we assume such a set exists and construct one rather than proving it \cite{nonfactoid_taxonomy2022}
      \item the answer structure is determined by the semantic context of the question and can be inferred from the question alone --- establishing this question-to-structure mapping in general is beyond the scope of this paper \cite{liroth2002questions}
    \end{itemize}
  \item \textbf{Scope:} we test the framework on multiple choice QA only, as a proof of concept; we hypothesize --- but do not claim to demonstrate --- that both the taxonomy and the approach extend to the general question space, which we leave to future work
  \item Evaluation: compare baseline small LLM vs.\ scaffolded small LLM across multiple model sizes and configurations, including base and instruction-tuned variants
\end{itemize}

% ============================================================
\section{Literature Review}
% ============================================================

\subsection{Large Language Models}

\begin{itemize}
  \item Large language models are deep neural networks built on the transformer architecture, trained on massive text corpora \cite{vaswani2017attention, llmoverview2023}
    \begin{itemize}
      \item Self-attention allows every token to attend to every other token in the sequence --- this is what enables the model to condition on provided context (including prompts and scaffolds) \cite{vaswani2017attention}
      \item ``Large'' refers to parameter count: the learned weights that encode statistical patterns from training, typically hundreds of millions to hundreds of billions
      \item The model learns by adjusting parameters to minimize prediction error across the training corpus --- what it encodes is bounded by its capacity to fit that corpus
    \end{itemize}
  \item Their function is probabilistic token prediction: given a sequence, the model outputs a probability distribution over the next token --- they are prediction engines, not knowledge databases \cite{brown2020gpt3, llmoverview2023}
    \begin{itemize}
      \item The final layer produces a probability distribution over the full vocabulary --- every possible next token gets a likelihood
      \item The model does not ``decide'' an answer --- it assigns likelihoods to all possible continuations based on patterns absorbed during training
      \item Output is inherently stochastic: the same input can produce different outputs, and nothing in the architecture guarantees correctness or consistency
    \end{itemize}
  \item The pre-training objective determines what the model learns to do --- and separates the two architectures relevant to this paper \cite{devlin2019bert, brown2020gpt3}
    \begin{itemize}
      \item Decoders (GPT-style): trained via next-token prediction --- see only left context, learn to generate \cite{brown2020gpt3}
      \item Encoders (BERT-style): trained via masked language modeling --- see full context bidirectionally, learn to represent \cite{devlin2019bert}
      \item This paper uses both: a decoder as the generator (small LLM) and an encoder as the classifier (answer structure type prediction)
    \end{itemize}
  \item Decoder LLMs generate via autoregressive decoding: one token at a time, each conditioned only on what came before --- no lookahead, no backtracking, no revision \cite{brown2020gpt3, exposure_bias_ilp2022}
    \begin{itemize}
      \item Once a token is emitted, it becomes part of the context for all subsequent tokens --- early commitments propagate forward
      \item If the model begins generating in the wrong direction, it has limited mechanism for self-correction under standard decoding
      \item This is confirmed architecture for all open models\footnote{Closed-source models may employ internal techniques (e.g.\ beam search, speculative decoding, hidden chain-of-thought) that partially mitigate this. However, the small open models this paper targets are confirmed autoregressive decoders \cite{llama3_2024, liesenfeld2024closed}.}
      \item This is the core vulnerability of autoregressive generation: early token commitments constrain all subsequent generation with no opportunity to revise
    \end{itemize}
  \item \textbf{Key implication:} the model is a probability engine --- anything that improves output quality does so by shaping the probability distribution over next tokens \cite{prompt_control_theory2023, llmreasoning_stochastic2024}
\end{itemize}

\subsection{Scale and Structure}

\begin{itemize}
  \item Scale, as used in the literature, refers primarily to the number of learnable parameters in the model. It is the dominant variable in scaling law research and the conventional predictor of general performance \cite{kaplan2020scaling}
  \item Larger models consistently outperform smaller ones on general tasks. This is well-established in the literature, formalized as scaling laws, and confirmed in industry deployment \cite{kaplan2020scaling, hoffmann2022chinchilla}
  \item The conventional explanation for this advantage is capacity \cite{kaplan2020scaling}. More parameters allow more connections between concepts to be encoded, and with more connections, more ideas can be held together during generation\footnote{Other factors contribute to the performance gap between model sizes, including training data distribution, optimization procedure, and architectural improvements. Training data can be identical across model sizes, so data alone does not explain the gap. Architectural improvements may also contribute, but since many frontier models are closed-source, their internal mechanisms cannot be verified externally \cite{liesenfeld2024closed}. The capacity argument remains central in the literature and is the focus of this section.} \cite{llmoverview2023}
  \item The jump from GPT-2 (1.5B parameters) to GPT-3 (175B parameters) demonstrated this dramatically: GPT-3 performed few-shot tasks that GPT-2 could not perform at all, without any change to architecture or training objective \cite{brown2020gpt3}
  \item Despite this, even at large scale, models produce unexpected output on tasks requiring multi-step logic or structured responses. This is well-documented and systematic \cite{math_reasoning_limits2024, lrm_complexity2025}
  \item More parameters does not mean every valid connection or context is encoded. Scale increases coverage but cannot guarantee completeness for an arbitrary task \cite{kaplan2020scaling, costofdownscaling2023}
  \item If the gap were purely a matter of capacity, external context should not close it. Yet on specific tasks and domains, small models match or exceed large ones when given appropriate context or specialization. The gap is conditional, not absolute \cite{smalllmsready2024, modelsizematters2024}
  \item Small models are not architecturally distinct from large models. They share the same training objective, generation mechanism, and target distributions. What differs is the dimensionality available to represent those distributions; scale determines the resolution of the approximation, not its kind \cite{godey2024softmax, costofdownscaling2023}
  % NOTE: REVISIT — review wording on next pass. This claim establishes architectural equivalence so §2.3 can reference it. Verify godey2024softmax supports "same distributions, fewer dimensions" framing specifically.
  \item We argue that part of the large-model advantage is structural. Larger models have the capacity to encode patterns for how answers should be structured, not just what they should contain. They have seen enough training examples to statistically represent answer structure within their parameters
  % NOTE: REVISIT — need supporting citations for the claim that large models encode answer structure specifically. Current argument is author synthesis only. Find papers showing format/structure sensitivity in large models or evidence of structural priors in pretrained weights.
  \item Because this structure is statistically encoded, it is never deterministic. There is no guarantee the correct structural pattern activates for a given input, even in large models \cite{math_reasoning_limits2024, lrm_complexity2025}
  \item External structure in the input partially compensates for what is not reliably encoded in parameters. Chain of Thought prompting imposes step-by-step structure in the input and measurably improves reasoning performance \cite{wei2022cot}
  \item Brevity constraints demonstrate it more directly: when output form is constrained, small models outperform large ones on mathematical reasoning and scientific knowledge by 7--16 percentage points \cite{brevity_constraints2025}
  \item Neither CoT nor brevity constraints are matched to the specific question. The gap remains: no existing technique classifies the question and provides answer structure matched to it
\end{itemize}

\subsection{Unconstrained Prediction}

\begin{itemize}
  \item What is commonly called ``hallucination'' is output that does not match the expected or correct answer for a given task \cite{ji2023hallucination, hallucination_unified2024}
  \item However, at the architectural level, the model's generation process is autoregressive next-token prediction: regardless of internal complexity, every forward pass produces a probability distribution over the next token (see \S2.1). This operation is the same whether the output is task-correct or task-incorrect. We argue that any token the model produces is a statistically probable continuation given its parameters and context, and that there is no separate failure mode that produces hallucinations \cite{brown2020gpt3, llmreasoning_stochastic2024}
  \item The model's parameters determine the baseline distribution over possible outputs. In small models, this baseline is less precise: as established in \S2.2, they approximate the same target distributions as large models but with fewer dimensions. The mechanistic basis for this is a dimensional bottleneck; the target distribution is high-rank, and when a model's hidden dimension falls below what is required to express it, the approximation degrades \cite{godey2024softmax, costofdownscaling2023}
  % NOTE: REVISIT — review wording on next pass. Depends on §2.2 architectural equivalence addition. Verify godey2024softmax supports this framing.
  \item This approximation compounds during autoregressive generation: each token is produced from an imprecise representation of what should come next, and that imprecision propagates forward through the sequence with limited mechanism for self-correction under standard decoding \cite{lpsr2026, exposure_bias_ilp2022}
  \item We argue that this compounding cannot be eliminated from autoregressive systems, but its scope can be reduced by narrowing the space of probable outputs before generation begins. Empirical evidence supports this: constraining output form alone inverts model-size performance rankings on reasoning tasks (see \S2.2) \cite{prompt_control_theory2023, brevity_constraints2025}
  \item Existing approaches reduce the prediction space through two mechanisms. Modifying parameters (training, fine-tuning) reshapes the baseline distribution itself, making some outputs permanently more or less likely. Enriching context (retrieved knowledge, reasoning processes, topic-similar examples) constrains the distribution at inference time by providing additional signal the model conditions on (see \S2.5, \S2.6)
  \item Answer structure (defined in \S2.8) reduces the prediction space by providing the output form before generation begins, so that the model's remaining task is to populate it with relevant content. This reduces the probability mass on structurally incorrect continuations regardless of the model's parametric coverage \cite{llmreasoning_stochastic2024}
  % NOTE: May need workshopping later.
\end{itemize}

\subsection{Small Language Models}

\begin{itemize}
  \item Small LLMs (roughly 1B--7B parameters) share the same autoregressive decoder architecture as large models but with fewer parameters \cite{slm_survey2024}
  \item They are accessible to individual researchers and organizations without large compute budgets \cite{slm_survey2024, smalllmsready2024}
  \item Small models retain core linguistic competence --- language understanding does not degrade proportionally with parameter count \cite{linguistic_competence_scale2025, emergent_reducedscale2024, costofdownscaling2023}
  \item What degrades is parametric fact recall and associative depth, not the ability to process and follow in-context information \cite{costofdownscaling2023}
  \item Domain specialization: small models can be trained effectively on specific domains, matching expert-level performance \cite{domain_expert_smalllm2024}
  \item But specialization comes at a cost --- a model trained for one domain cannot easily transfer to another without catastrophic forgetting \cite{domain_expert_smalllm2024, catastrophic_forgetting2023}
  \item On general tasks, small models still underperform their larger counterparts \cite{modelsizematters2024, kaplan2020scaling}
  \item \textbf{Key implication:} the deficit in small models is not linguistic competence but organizational structure --- they understand language but lack the capacity to independently organize what they retain into correct answers for arbitrary questions
\end{itemize}

\subsection{Fine-Tuning and Parameter-Efficient Adaptation}

\begin{itemize}
  \item Full fine-tuning: update all model weights on task-specific data --- effective but computationally expensive and requires large labeled datasets \needpaper{} (general fine-tuning survey)
  \item Catastrophic forgetting: fine-tuning on a new task degrades performance on previously learned tasks \cite{catastrophic_forgetting2023}
  \item Parameter-Efficient Fine-Tuning (PEFT): update only a small subset of parameters to reduce cost
    \begin{itemize}
      \item LoRA: inject low-rank trainable matrices into frozen layers \needpaper{} Hu et al.\ 2022 ``LoRA''
      \item Adapters: insert small trainable modules between frozen layers \needpaper{} Houlsby et al.\ 2019
      \item Prefix tuning: prepend trainable continuous vectors to the input \needpaper{} Li \& Liang 2021
    \end{itemize}
  \item PEFT reduces compute and storage costs but still requires task-specific training data and a training loop per task \needpaper{} (PEFT survey)
  \item Instruction tuning: fine-tune on instruction-following data to improve zero-shot generalization \cite{instructiontuning2023}
    \begin{itemize}
      \item Improves usability but still requires a training phase; cannot be applied to arbitrary untrained models at inference time
    \end{itemize}
  \item \textbf{Shared limitation:} all weight-modification approaches are task-specific --- they improve the model for one distribution but require retraining for each new task or domain, and risk forgetting previous capabilities
  \item \textbf{Distinction from this paper:} our framework requires \textit{no weight updates} --- it operates entirely at inference time and works with any small LLM as-is, whether base or instruction-tuned
\end{itemize}

\subsection{Prompt Engineering and In-Context Learning}

\begin{itemize}
  \item Prompt engineering: shaping model output at inference time via input construction, without modifying weights \cite{brown2020gpt3}
  \item \textbf{In-Context Learning (ICL):} models learn tasks from examples provided in the prompt, without weight updates \cite{brown2020gpt3}
    \begin{itemize}
      \item ICL can be understood as implicit Bayesian inference --- few examples sharply reduce the model's uncertainty over the latent task \cite{xie2022icl}
      \item Critically, ICL ability survives aggressive down-scaling even as parametric fact recall degrades \cite{costofdownscaling2023}
      \item This is the mechanism our framework exploits: the scaffold works because small models can still \textit{use} provided structure via ICL
    \end{itemize}
  \item \textbf{Chain-of-Thought (CoT):} elicit step-by-step reasoning before the final answer \cite{wei2022cot}
    \begin{itemize}
      \item Works well on multi-step reasoning for large models
      \item Unreliable on small models; the underlying mechanism is debated \cite{cot_factors2024, cot_limitations2025}
      \item Recent evidence: CoT may mimic reasoning patterns without performing genuine abstract reasoning \cite{cot_mirage2025}
      % NOTE: Argument to develop here — CoT does not provide external structure; it lets the model generate its own reasoning, which is just autoregressive generation biased toward reasoning-sounding language. Gains come from additional context each iteration provides, not from genuine structural guidance. Frame as "we argue" and cite cot_mirage2025 + cot_limitations2025 as supporting evidence. Be careful not to overstate — CoT does measurably improve large model performance; the argument is about mechanism and about failure on small models.
    \end{itemize}
  \item \textbf{CoT extensions:}
    \begin{itemize}
      \item Tree of Thoughts: explore multiple reasoning paths as a tree search \cite{yao2023tot}
      \item Self-Consistency: sample multiple CoT chains, select most consistent answer \cite{wang2023selfconsistency}
      \item Limitation: computationally expensive; still knowledge-dependent
    \end{itemize}
  \item \textbf{Few-shot / many-shot prompting:} provide example input-output pairs to condition generation \cite{brown2020gpt3, manyshot2024}
    \begin{itemize}
      \item Examples are selected by semantic similarity to the query topic, \textbf{not} by answer structure type --- this is the gap our work fills
    \end{itemize}
  \item \textbf{Mixture of Experts (MoE) as a related routing concept:} MoE divides a model into specialized sub-networks routed by a learned gate \cite{shazeer2017moe, fedus2022switchtransformer}. Conceptually similar to our classifier-based routing, but MoE operates at training time and requires massive infrastructure; our classifier operates at inference time with no retraining of the generator \cite{moe_survey2024}
  \item \textbf{Summary of gap:} all existing prompt engineering techniques construct prompts based on content or reasoning paths. None classify the \textit{structural type of answer required} before generating. This is the novel contribution of the present framework.
\end{itemize}

\subsection{Retrieval-Augmented Generation}

\begin{itemize}
  \item RAG: augment a model's parametric knowledge by retrieving relevant documents at inference time, then conditioning generation on the retrieved context \cite{lewis2020rag}
  \item RAG variants:
    \begin{itemize}
      \item Naive, Advanced, Modular RAG (architectural taxonomy) \cite{rag_survey2023}
      \item Speculative RAG: small model drafts, large model verifies \cite{speculative_rag2024}
      \item Question-decomposition RAG: decompose the question into sub-questions for retrieval \cite{question_decomp_rag2025}
    \end{itemize}
  \item \textbf{Critical differentiation from the present work:}
    \begin{itemize}
      \item RAG addresses \textit{what information} the model receives (knowledge augmentation)
      \item Question-decomposition RAG addresses how the question is broken into sub-questions (content decomposition)
      \item This paper addresses \textit{what structural type} the answer requires (answer structure classification)
      \item These are three distinct problems --- our approach is orthogonal to RAG and could be combined with it
    \end{itemize}
  \item The few-shot scaffold in our framework is conceptually similar to RAG's retrieved context: both prepend information to the prompt at inference time. Key difference: RAG retrieves factual content; we retrieve structural templates.
\end{itemize}

\subsection{Question Classification in QA}

\begin{itemize}
  \item Classical QA research has long recognized question type classification as a necessary preprocessing step in information retrieval systems \cite{nonfactoid_taxonomy2022} --- \needpaper{} TREC question classification (Li \& Roth 2002)
  \item \needpaper{} Cognitive science / human reasoning taxonomy (Bloom's Taxonomy or Kahneman dual-process theory)
  \item \textbf{Cognitive science grounding:} expert human problem-solvers categorize a problem type before attempting a solution (dual-process theory: System 2 reasoning begins with problem classification) \cite{pnas_analogy2010} --- \needpaper{} Kahneman dual-process review
  \item Different question types require structurally different reasoning processes:
    \begin{itemize}
      \item \textbf{Factual:} retrieve and state a known fact
      \item \textbf{Analogical:} map relational structure from one domain to another \cite{analogical_reasoning2022}
      \item \textbf{Multi-step:} chain intermediate conclusions
      \item \textbf{Opinion / non-factoid:} construct a structured argument \cite{nonfactoid_taxonomy2022}
    \end{itemize}
  \item Prior work classifies \textit{question types} --- but does not define or classify the \textit{structure that a valid answer must take} for each type \cite{liroth2002questions, nonfactoid_taxonomy2022}
\end{itemize}

\subsection{Answer Structure Types: A Taxonomy Hypothesis}

\begin{itemize}
  \item Not all questions require the same type of answer --- a factual recall question requires a different response structure than an analogical reasoning or multi-step problem
  \item \textbf{Hypothesis of this paper:} these structural types are finite, enumerable, and classifiable from question text alone (no answer required); a BERT-style encoder can be trained to classify them reliably \cite{devlin2019bert}
  \item An answer structure defines the \textit{dimensions/components} a valid answer must contain --- not the answer content itself
  \item \textbf{Implication:} if structure classification is natural for humans and classical QA systems, a trained encoder can replicate it --- and the resulting scaffold can guide a small LLM to the correct answer type
\end{itemize}

% ============================================================
\section{Architecture}
% ============================================================

\begin{itemize}
  \item \textbf{Overview:} two-stage inference-time pipeline, no weight updates required
    \begin{itemize}
      \item Stage 1 --- Classification: encoder model classifies question $\rightarrow$ answer structure type
      \item Stage 2 --- Generation: small LLM generates answer conditioned on retrieved scaffold
    \end{itemize}
  \item \textbf{Stage 1 --- The Classifier}
    \begin{itemize}
      \item Architecture: BERT-style bidirectional encoder \cite{devlin2019bert}
      \item Input: raw question text
      \item Output: answer structure type label (from the taxonomy)
      \item Training: fine-tuned on labeled question-type dataset \cite{datasets_survey2024}
      \item Why an encoder: bidirectional attention gives richer question representations than a left-to-right decoder; classification is a natural fit \cite{devlin2019bert, liu2019roberta}
    \end{itemize}
  \item \textbf{The Taxonomy}
    \begin{itemize}
      \item Define the set of answer structure types used in this paper (e.g.\ factual-recall, analogical, multi-step reasoning, multiple-choice selection)
      \item Discuss how the taxonomy was derived
      \item Justify that the taxonomy is principled, not arbitrary \cite{nonfactoid_taxonomy2022}
    \end{itemize}
  \item \textbf{Stage 2 --- Scaffold Retrieval}
    \begin{itemize}
      \item Each taxonomy class has a corresponding pre-constructed few-shot prompt scaffold
      \item Scaffold structure: [answer type description] + [N worked examples] + [question]
      \item The scaffold narrows the probability space before generation begins \cite{brown2020gpt3, manyshot2024}
    \end{itemize}
  \item \textbf{Stage 2 --- The Generator}
    \begin{itemize}
      \item Any small, English-understanding LLM (baseline or instruction-tuned)
      \item No fine-tuning required
      \item The scaffold does the structural work; the model supplies language understanding \cite{smalllmsready2024}
    \end{itemize}
  \item \textbf{Why this works (theoretical argument):}
    \begin{itemize}
      \item LLMs are probability engines; the scaffold shifts the prior toward valid response structures
      \item This is analogous to how few-shot prompting works, but at the structural level rather than the topical level \cite{llmreasoning_stochastic2024, brown2020gpt3}
    \end{itemize}
\end{itemize}

% ============================================================
\section{Methodology}
% ============================================================

\begin{itemize}
  \item \textbf{Research questions:}
    \begin{itemize}
      \item RQ1: Does structural scaffolding improve small LLM accuracy on multiple choice QA?
      \item RQ2: Is the improvement consistent across different model sizes?
      \item RQ3: Does the classifier reliably identify answer structure types?
    \end{itemize}
  \item \textbf{Dataset selection}
    \begin{itemize}
      \item Multiple choice QA benchmarks (e.g.\ ARC, MMLU, HellaSwag) --- \needpaper{} ARC and MMLU dataset citations
      \item Justify: clear correctness criteria, well-established benchmarks with diverse question types
    \end{itemize}
  \item \textbf{Models evaluated}
    \begin{itemize}
      \item Range of small LLMs across parameter counts (e.g.\ 125M, 350M, 1B, 3B)
      \item Both base and instruction-tuned variants
      \item Specific models TBD based on availability
    \end{itemize}
  \item \textbf{Experimental conditions}
    \begin{itemize}
      \item \textbf{Baseline:} small LLM, zero-shot, no scaffold
      \item \textbf{Treatment:} small LLM + classifier-generated structural scaffold
      \item \textbf{Ablation:} small LLM + random scaffold (wrong structure type) --- tests whether classification quality matters or scaffolding volume alone drives improvement
      \item \textbf{Oracle upper bound:} small LLM + correct scaffold manually assigned --- sets the ceiling for a perfect classifier
    \end{itemize}
  \item \textbf{Classifier training}
    \begin{itemize}
      \item Training data: labeled question-type examples
      \item Architecture: fine-tuned BERT / RoBERTa \cite{devlin2019bert, liu2019roberta}
      \item Evaluation: accuracy of structure type prediction
    \end{itemize}
  \item \textbf{Evaluation metrics}
    \begin{itemize}
      \item Primary: accuracy on multiple choice questions
      \item Secondary: classifier accuracy on answer structure type prediction
    \end{itemize}
\end{itemize}

% ============================================================
\section{Results}
% ============================================================

\begin{itemize}
  \item \textbf{Classifier performance:} accuracy of answer structure type classification --- \textbf{TBD}
  \item \textbf{Main result:} scaffolded vs.\ baseline accuracy per model size --- present as table: model size $\times$ condition $\times$ accuracy --- \textbf{TBD}
  \item \textbf{Ablation result:} random scaffold vs.\ correct scaffold --- tests whether structure classification quality drives the improvement --- \textbf{TBD}
  \item \textbf{Oracle upper bound:} how much headroom remains if classification were perfect --- \textbf{TBD}
  \item \textbf{Qualitative examples:} cases where scaffolding helped and where it did not --- \textbf{TBD}
\end{itemize}

% ============================================================
\section{Discussion}
% ============================================================

\begin{itemize}
  \item \textbf{Main result interpretation:} does structural scaffolding close the gap between small and large models on these tasks?
  \item \textbf{Scaling behavior:} does improvement scale with model size or is it consistent across sizes?
    \begin{itemize}
      \item If consistent: supports the claim that language understanding is sufficient and the bottleneck is structural guidance, not capacity
    \end{itemize}
  \item \textbf{Ablation interpretation:} if random scaffold underperforms correct scaffold, this validates that answer structure classification matters, not just scaffolding volume
  \item \textbf{Limitations:}
    \begin{itemize}
      \item Experimental scope limited to multiple choice QA
      \item Taxonomy is hand-designed; broader coverage is future work
      \item Results may not generalize to open-ended generation tasks
      \item Classifier performance is a ceiling on overall system performance
    \end{itemize}
  \item \textbf{Connection back to problem statement:} results support the claim that unguided generation is the bottleneck, not knowledge capacity
  \item \textbf{Comparison to related work:} how do results compare to CoT-prompted and RAG-augmented small models? \cite{wei2022cot, lewis2020rag}
\end{itemize}

% ============================================================
\section{Limitations}
% ============================================================

\begin{itemize}
  \item \textbf{TBD}
\end{itemize}

% ============================================================
\section{Conclusion}
% ============================================================

\begin{itemize}
  \item \textbf{Restate the problem:} small LLMs fail not from lack of language understanding but from unguided generation --- no structural entry point before generating
  \item \textbf{Restate the contribution:} a two-component inference-time framework (encoder classifier + few-shot scaffold) that requires no retraining of the generator
  \item \textbf{Restate the central finding:} structural scaffolding consistently improves small LLM performance on multiple choice QA across model sizes --- \textbf{TBD with results}
  \item \textbf{Restate the taxonomy hypothesis:} a finite set of answer structure types is sufficient to cover the experimental task space; broader generalization is a direction for future work
  \item \textbf{Future work:}
    \begin{itemize}
      \item Extend the taxonomy to open-ended question types
      \item Apply the framework to other small model families
      \item Combine with RAG for knowledge-augmented + structure-guided generation
      \item Explore whether the framework reduces hallucinations specifically \cite{ji2023hallucination}
      \item Investigate whether the classifier can be made domain-agnostic
    \end{itemize}
\end{itemize}

\bibliographystyle{alpha}
\bibliography{sample}

\end{document}
